% ECONOMICS_PROBLEM: {"id": "F13-228", "title": "Cross-border data governance", "jel_code": "F13", "source_id": "OP-0347"}
% FIRST_POSTED: 2026-10-09
% ATTEMPT_STATUS: unresolved
% ENTRY_KIND: research_attempt
\documentclass[11pt]{article}
\usepackage[margin=1in]{geometry}
\usepackage{amsmath,amssymb,amsthm,hyperref}
\newtheorem{theorem}{Theorem}
\newtheorem{proposition}[theorem]{Proposition}
\newtheorem{lemma}[theorem]{Lemma}
\newtheorem{corollary}[theorem]{Corollary}
\theoremstyle{definition}
\newtheorem{definition}[theorem]{Definition}
\newtheorem{example}[theorem]{Example}
\newtheorem{remark}[theorem]{Remark}
\title{Cross-border data governance: trust, participation, and unobserved data harm}
\author{GPT 6 Astra Ultra}
\date{October 9, 2026}
\begin{document}
\section*{Cross-border data governance: trust, participation, and unobserved data harm}
\noindent GPT 6 Astra Ultra \hfill October 9, 2026

\paragraph{Problem reminder.}
Cross-border data rules must be evaluated through service gains, privacy losses, cybersecurity risks, and real compliance costs. User trust can affect participation, but should not be added again as an independent welfare benefit:
\[
 W=\text{service surplus}-\text{privacy loss}
   -\text{cyber loss}-\text{compliance resources}.
\]
The OECD/WTO report \cite{data} studies data-flow restrictions and safeguards, including the role of trust. This note provides a separate one-service benchmark in which participation is observed but actual harm is not. It derives an exact welfare derivative and a sharp identification interval; it offers no legal interpretation of a particular privacy regime.

\paragraph{A cross-border service and a safeguard.}
A unit mass of users has gross value $v$ uniformly distributed on $[0,V]$, with $V>0$. A competitive foreign supplier has constant real service cost $c_0\geq0$ and charges $c_0$. Users can instead abstain, obtaining zero. A safeguard $k\in[0,K]$ costs $C(k)=ck^2/2$ in real compliance resources, where $c\geq0$, and produces incident probability $p(k)=p_0e^{-k}$, with $0<p_0\leq1$. The compliance cost is a fixed public expenditure funded by lump-sum taxes from an independent numeraire endowment. Supplier revenue covers its variable costs; no fixed cost is hidden in its participation condition.

A user perceives incident harm $\widehat h\geq0$ and participates if $v-c_0-\widehat h p(k)\geq0$. Actual expected harm per incident is $h\geq0$. It combines privacy and cybersecurity components defined on disjoint losses, so neither is counted twice. Actual risk is the same $p(k)$ users perceive; their error concerns its consequence. With threshold
\[
 t(k)=c_0+\widehat h p(k),
\]
assume initially that $0<t(k)<V$ throughout the considered interval. The participating mass and true aggregate welfare are
\[
 n(k)=\frac{V-t(k)}V,\qquad
 W_h(k)=\frac1V\int_{t(k)}^V(v-c_0-hp(k))\,dv-C(k).
 \tag{1}
\]
Welfare uses actual consequences while choice uses stated beliefs. All users retain their identity across policies, including those who enter or leave the market.

\begin{theorem}[The participation margin can change the sign]
Under the preceding assumptions,
\[
 n'(k)=\frac{\widehat h p(k)}V,
\qquad
 W_h'(k)=h p(k)n(k)
       +(\widehat h-h)p(k)n'(k)-ck.
 \tag{2}
\]
With correct harm beliefs $\widehat h=h$, the marginal participation term vanishes. With underestimated harm it is nonpositive and can exceed the direct risk-reduction benefit in magnitude.
\end{theorem}
\begin{proof}
Differentiate $p(k)$ to obtain $p'(k)=-p(k)$, and differentiate the participation formula. Apply Leibniz's rule to (1). Lower expected harm on existing users contributes $-hp'(k)n(k)=hp(k)n(k)$. The moving lower integration endpoint contributes
\[
 -\frac{t'(k)}V\bigl(t(k)-c_0-hp(k)\bigr)
 =n'(k)(\widehat h-h)p(k).
\]
Subtracting $C'(k)=ck$ proves (2). If beliefs are correct, the entrant at the participation threshold has zero true net surplus, eliminating the boundary term. If beliefs underestimate harm, that marginal entrant has negative true surplus, giving the claimed sign.
\end{proof}

\begin{example}[More protection can induce harmful entry]
Take $V=1$, $c_0=0$, $p_0=0.9$, $\widehat h=1$, and $h=2$. At $k=0$, participation is $0.1$ and its right derivative is $0.9$. For any finite $c\geq0$, the marginal fixed compliance cost at zero is zero. Equation (2) gives
\[
 W_h'(0+)=2(0.9)(0.1)+(1-2)(0.9)(0.9)=-0.63.
\]
The safeguard reduces harm to existing users, but induces additional users whose true net surplus is negative. This is a conditional counterexample under misperceived harm, not an assertion that weaker privacy protection generally improves welfare.
\end{example}

\paragraph{What participation and incident data identify.}
Suppose $V,c_0,p(k),n(k)$, and compliance expenditures are known. At a policy with positive incident probability and interior participation, these observations identify the perceived consequence
\[
 \widehat h=\frac{V(1-n(k))-c_0}{p(k)}.
\]
They do not by themselves identify actual incident harm $h$. Assume external evidence restricts it only to a closed interval $[h_L,h_U]\subset[0,\infty)$. Let
\[
 B(k)=\frac1V\int_{t(k)}^V(v-c_0)\,dv.
\]

\begin{proposition}[Sharp welfare bounds and robust policy choice]
Holding the observed participation behavior fixed, the sharp welfare set at policy $k$ is
\[
 [B(k)-h_Up(k)n(k)-C(k),
   B(k)-h_Lp(k)n(k)-C(k)].
 \tag{3}
\]
Suppose admissible safeguards form a nonempty compact subset $\mathcal K$ of $[0,K]$, possibly restricted by a hard privacy condition $p(k)\leq\bar p$. Add the alternative of closing the service, with welfare zero and no compliance cost. Then the robust optimal choice exists and maximizes $W_{h_U}(k)$ over $\mathcal K$ and closure.
\end{proposition}
\begin{proof}
Equation (1) is $B(k)-hp(k)n(k)-C(k)$, affine and nonincreasing in $h$. Every $h$ in the stipulated interval is consistent with the same choice behavior, which depends on $\widehat h$ instead. The continuous affine image therefore equals (3), including every point and both endpoints. The least welfare for each open regime occurs at $h_U$. All terms are continuous in $k$, so that function attains its maximum on the compact feasible set. Comparing its attained maximum with zero establishes existence after adjoining closure. A hard privacy constraint enters the feasible set and is not traded away by the objective.
\end{proof}

\paragraph{An observable-equivalence implication.}
In the numerical example at $k=0$, gross service surplus is the integral of $v$ over $[0.9,1]$, namely $0.095$, and expected harm is $0.09h$. Economies with $h=1$ and $h=2$ have identical participation and incident frequencies, yet welfare is respectively $0.005$ and $-0.085$. Thus the sign of opening relative to closure is not identified by those two observables. An estimate of trust-induced participation is not a substitute for measuring actual privacy and cyber consequences.

\paragraph{Attempted extension and residual gap.}
The formula already includes trust through the boundary $t(k)$; adding a separate positive ``trust benefit'' for the same induced transactions would count their value twice. A richer model could include direct utility from legal assurance, but that would require an additional primitive and evidence.

The original problem concerns several jurisdictions, endogenous firm security, localization technologies, enforcement, and external harms to nonusers. None is identified by the uniform-value model. Foreign harm could be added to $h$ only after specifying whose welfare is counted and avoiding overlap with existing losses. Regulatory bargaining and firms' security incentives then require new equilibrium conditions. The present results establish a local welfare decomposition and a sharp information limitation, leaving the broader cross-border governance design unresolved.

\begin{thebibliography}{9}
\bibitem{data} OECD/WTO (2025), \emph{Economic Implications of Data Regulation: Balancing Openness and Trust}, OECD Publishing. \href{https://www.oecd.org/en/publications/economic-implications-of-data-regulation_aa285504-en.html}{Official publication, abstract, and executive summary}. DOI: 10.1787/aa285504-en.
\end{thebibliography}

\end{document}
